\documentclass[12pt,a4paper,oneside]{article}
\usepackage[brazil]{babel}
\usepackage[utf8]{inputenc}
\usepackage{olymp}

\setcounter{problem}{0}

\begin{document}

\begin{problem}{Sequência de collatz}{collatz}{2}

A sequência de Collatz (conhecida como $3n+1$) usa uma regra matemática simples a partir de número inteiro positivo $n$:

\begin{itemize}
\item Se $n$ for par: divida por $2$, isto é, $n/2$;
\item Se $n$ for ímpar: multiplique por $3$ e some $1$, isto é, $3n+1$;
\item Repita o processo com o novo número até que ele se torne $1$.
\end{itemize}

Exemplo: Se iniciar com $n=6$, temos a sequência $6 \rightarrow 3 \rightarrow 10 \rightarrow 5 \rightarrow 16 \rightarrow 8 \rightarrow 4 \rightarrow 2 \rightarrow 1$.

Sua tarefa aqui é descobrir a quantidade de passos para um dado valor inicial $n$ chegar em $1$.
Por exemplo, para $n=6$, são necessários 8 passos.

%\Notes
%
%Observações, se tiver.

\InputFile

A entrada contém um único número natural $N$. Restrições: $2 \leq N \leq 200$.

\OutputFile

Seu programa deve informar o número de passos da sequência de Collatz para o valor $N$.

%\Notes

\Example

\begin{example}
\exmp{
6
}{
8
}%
\end{example}

\begin{example}
\exmp{
4
}{
2
}%
\end{example}

\begin{example}
\exmp{
27
}{
111
}%
\end{example}

%Comentários sobre o exemplo, se necessário.
\Notes
Curiosidade: a conjectura de Collatz, proposta pelo matemático Lothar Collatz em 1937, é um dos mais famosos problemas não resolvidos da Matemática.
Ela pergunta se a sequência de Collatz eventualmente chega em $1$ para qualquer número natural $n$. Já foi mostrado que vale para todo $n < 2.36 \times 10^{31}$, mas até hoje não se sabe se é verdadeira para qualquer $n$.

\end{problem}

\end{document}
